\section{Motivación de RLP: del texto observacional al aprendizaje mediante pensamiento}

RLP ya no solo ajusta el orden de los datos, sino que introduce una acción pensable (thought action) muestreable antes de cada token seleccionado. La idea central es: el modelo genera primero varias trazas de razonamiento interno y luego predice el siguiente token real; si cierta thought aumenta la probabilidad de ese token, recibe una recompensa positiva. De este modo, el texto de preentrenamiento estándar proporciona simultáneamente el token objetivo y una referencia de recompensa auto-supervisada.

\subsection{La historia de dos aprendices}

Leo primero construye un puente rudimentario pero funcional para un vehículo pequeño, luego lo mejora a partir de la retroalimentación; Bolt lee todo el material sobre puentes y crea modelos hermosos que no cumplen las restricciones. La historia destaca la diferencia entre aprender haciendo (learning by doing) y aprendiendo observando (learning by observing). La predicción del siguiente token es hábil imitando patrones estadísticos a partir de textos ya escritos, mientras que RLP intenta otorgar al modelo una oportunidad de exploración para «proponer primero una explicación y luego verificarla con el token real».

\lecturefigure{slide-023.jpg}{Diferencia entre dos aprendices: ensayo y error práctico versus observación de ejemplos}{Grabación oficial de Stanford, 00:28:14--00:29:37}

\begin{warningbox}{Thought no es una explicación automáticamente confiable}
Las cadenas de pensamiento (chain-of-thought) generadas por el modelo pueden contener racionalizaciones a posteriori, hechos erróneos o atajos incomprensibles. RLP solo verifica si la thought aumenta la probabilidad de predicción del siguiente token observado; no valida si coincide con el razonamiento humano, si es causalmente fiel o si se mantiene fuera de la distribución.
\end{warningbox}

\subsection{El problema del preentrenamiento estándar no es la falta de gradientes, sino la ausencia de acciones de exploración}

El NTP estándar ya proporciona gradientes densos, pero el espacio de acciones (action space) solo permite ajustes directos a la distribución del siguiente token; el razonamiento complejo suele aparecer explícitamente recién en SFT, RLHF o RLVR. La formulación del problema de RLP es: ¿es posible anticipar la exploración manteniendo al mismo tiempo la escalabilidad del flujo documental estándar? Esta motivación también responde al muro de datos: si los nuevos tokens se vuelven cada vez más costosos, hay que aumentar la señal de aprendizaje generada por cada token.

\lecturefigure{slide-024.jpg}{El preentrenamiento estándar retrasa la exploración explícita hasta el post-entrenamiento}{Grabación oficial de Stanford, 00:29:37--00:31:23}

El objetivo del NTP estándar es
\[
\mathcal{L}_{\mathrm{NTP}}(\theta)
=-\sum_{t=1}^{L}\log p_\theta(x_t\mid x_{<t}).
\]
Donde, $x_{<t}$ representa el contexto, $x_t$ es el siguiente token real en el corpus y $p_\theta$ es la distribución del modelo. Este objetivo aprende mucha estructura, pero no tiene una variable separada que represente «qué thought se exploró antes de generar la salida».

\subsection{Diferencia entre Vanilla y RLP en probabilidades condicionales}

En el ejemplo de la fotosíntesis, el modelo vanilla predice directamente `sunlight` a partir del contexto; el modelo RLP genera primero una thought que dice «este proceso depende de energía solar», luego condicionalmente al contexto y a esa thought predice el mismo token. La clave no es generar más texto, sino tratar la thought como una acción optimizable y comparar su ganancia relativa respecto a un baseline que solo usa el contexto.

\lecturefigure{slide-025.jpg}{Predicción vanilla del siguiente token versus predicción RLP condicionada al pensamiento}{Grabación oficial de Stanford, 00:31:23--00:32:34}

\[
p_\theta(x_t\mid x_{<t},c_t)
\quad\text{versus}\quad
p_\phi(x_t\mid x_{<t}),
\]
donde, $c_t$ es la thought muestreada antes de la posición $t$, $\theta$ son los parámetros del policy de thoughts o predictor razonado, y $\phi$ son los parámetros del baseline sin pensamiento. Ambos ven el mismo contexto y el mismo token objetivo real; la diferencia radica únicamente en si se incluye $c_t$.

\begin{importantbox}{Ciclo auto-supervisado de RLP}
El token real $x_t$ del corpus cumple simultáneamente dos funciones: es el objetivo para NTP y sirve como referencia para evaluar si la thought tiene utilidad predictiva. Por ello, RLP no requiere un verificador de respuestas manual, pero sí depende del siguiente token real; no recompensa de forma insupervisada cualquier texto que «parezca razonamiento».
\end{importantbox}

\subsection{Resumen de la sección}

RLP convierte la thought en una acción exploratoria antes de la predicción y verifica si dicha acción aporta información mediante el siguiente token real. La diferencia con NTP no es si existen señales de entrenamiento, sino si hay un camino de exploración explícito y recompensas relativas sobre la calidad del camino. A continuación se desglosa paso a paso el ciclo de entrenamiento completo.
